\begin{table*}[t]
\centering
\scriptsize
\setlength{\tabcolsep}{3.2pt}
\resizebox{\textwidth}{!}{%
\begin{tabular}{@{}llllcccccc@{}}
\toprule
\textbf{Study} & \textbf{Intervention} & \textbf{Adaptation} &
\textbf{Trajectory} & \textbf{Loss+task} & \textbf{Scratch/init} &
\textbf{Intact CPT} & \textbf{MMLU iface.} & \textbf{Closed-book} &
\textbf{Construction} \\
\midrule
Gromov et al. & deep-block removal & CPT & partial & \checkmark & --- & --- & --- & --- & depth/budget \\
Shortened LLaMA & depth pruning & CPT/LoRA & \checkmark & \checkmark & \checkmark & --- & --- & --- & ratio/method \\
Minitron & structured pruning & distillation/CPT & \checkmark & \checkmark & \checkmark & --- & --- & --- & iterative/shape \\
IteRABRe & iterative block removal & iterative recovery & \checkmark & \checkmark & --- & --- & --- & --- & iteration \\
PASER & pruned-model recovery & selected CPT data & partial & \checkmark & --- & --- & --- & --- & data policy \\
\textbf{This study} & prefix+fresh tail & CPT & \checkmark & \checkmark & op.\ point & 25k only & \checkmark & \checkmark & coupled 16L \\
\bottomrule
\end{tabular}%
}
\caption{\textbf{Nearest-work positioning.} Prior work already studies recovery
curves, loss--task gaps, initialization, and iterative retraining. The narrower
increment here is the \emph{combination} of an OLMo prefix+fresh-tail path with
an available short-horizon intact branch, same-shape operating points, two MMLU
interfaces, closed-book QA, and an explicitly coupled 16-layer construction
comparison. ``Partial'' denotes intermediate measurements that are not the sole
analysis object; ``op.\ point'' is not an initialization-only ablation.}
\label{tab:priorart}
\end{table*}
